\documentclass[10pt]{article}

\usepackage[utf8]{inputenc}

\usepackage[T1]{fontenc}

\usepackage{amsmath}

\usepackage{amsfonts}

\usepackage{amssymb}

\usepackage[version=4]{mhchem}

\usepackage{stmaryrd}

\usepackage{bbold}

\usepackage{graphicx}

\usepackage[export]{adjustbox}

\graphicspath{ {./images/} }



\begin{document}

ЧВЕТ - Квантовое число кварков и глюонов, характеризуюшее физически неразличимые унитарно эквивалентные состояния этих фундаментальных частиц. Ц. принимает для каждого типа (или аромата) кварков $q$ три значения и восемь значений для глюона $G$, то есть $q=\left\{q_{i}\right\}, i=1,2,3$, $G=\left\{G_{a}\right\}, a=1, \ldots, 8$. Ц. впервые введен в теорию элементарных частиц в 1965 в рамках динамич. составной кварковой модели адронов (см. [1]) как дополнительное квантовое число кварков, позволяющее разрешить проблему статистики кварков и дать последовательное объяснение закономерностей спектроскопии сильно взаимодействуюших частиц - мезонов и барионов (см. [2], [3]). При этом предполагалось, что волновые функции низколежаших мезонов $M$, состоящих из пары кварка и антикварка, и барионов $B$, состоящих из трех кварков, обладают максимальной симметрией по отношению к Ц. Кварков (п р и нцип б е с цветности адронов), то есть



\[

\begin{gathered}

M \sim \delta_{j}^{i} q_{i}(1) \bar{q}^{j}, \\

B \sim \in{ }^{i j k} q_{i}(1) q_{j}(2) q_{k}(3) .

\end{gathered}

\]



В настояшее время считается, что Ц. всех наблюдаемых состояний сильно взаимодействуюших частиц (адронов) равен нулю, тогда как обладаюшие отличным от нуля Ц. кварки и глюоны удерживаются в адронах субъядерными, так наз. хромоди намическими силам и и не могут наблюдаться в свободном состоянии.



Переносчиком хромодинамич. сил, как бы склеивающих кварки в адронах, являются соответствующие цветным глюонам безмассовые калибровочные векторные поля (см. [5]).



В современной калибровочной теории сильных взаимодействий (квантовой хромодинамике) Ц. ассоциируется с зарядом, принимающим, в отличие, напр., от электрич. заряда, значения в многомерном векторном пространстве (размерности, равной трем для кварков и восьми для глюонов). Как и электрич. заряд Ц. строго сохраняется, что обусловлено локальной природой цветовой симметрии в квантовой хромодинамике, описываемой неабелевой калибровочной группой $S U_{c}(3)$.



В ряде ранних работ в качестве альтернативы введению Ц. для устранения противоречия с принципом Паули в кварковой модели Гелл-Мана и Цвейга обсуждалась гипотеза парафермистатистики кварков ранга 3 (см. [6]). Оказывается, что максимальной группой симметрии в калибровочной теории парафермионов ранга 3 является ортогональная группа $S O(3)$, приводящая к трем калибровочным векторным бозонам (глюонам) вместо восьми, соответствуюшим группе $S U(3)$. Таким образом, гипотеза парафермистатистики кварков не является физически приемлемой альтернативой гипотезе цветных кварков, удовлетворяющих нормальной статистике Ферми - Дирака.



Лum.: [1] Бо голюбов Н. Н., С т рум и нский Б. В., Та в хе лидз е А. Н., К вопросу о составных моделях в теории элементарных частиц, Дубна, 1965; [2] Релятивистски инвариантные уравнения для составных частиц и формфакторы, Дубна, 1965; [3] Н а м бу Й., «Успехи физич. наук», 1978, т. 124, с. 147-69; [4] е го же, Кварки, М., 1984; [5] H a n M., N a mbu Y., "Phys. Rev.», 1965, v. 139, № 48, p. 1006-10; [6] N a m b u Y., A systematics of hadrons in subnuclear physics, Chicago, 1965; [7] Greenber O., "Phys. Rev. Lett.», 1964, v. 13, p. 598-602.



К. Г. Четыркин. ЧЕЙНЕЛЯ МЕТОД - см. Усреднения метод.



ЦЕЛАЯ РАЦИОНАЛЬНАЯ ФУНКЦИЯ - сМ. Целая функция, Рациональная функция, Особая точка аналитической функции.



ЦЕЛАЯ ТРАНСЧЕНДЕНТНАЯ ФУНКЧИЯ - см. Особая точка аналитической функции.



ЧЕЛАЯ ФУНКЦИЯ - функция, аналитическая во всей плоскости комплексного переменного (кроме, возможно, бесконечно удаленной точки). Она разлагается в степенной ряд



\[

f(z)=\sum_{k=0}^{\infty} a_{k} z^{k}, \quad a_{k}=\frac{f^{(k)}(0)}{k !}, \quad k \geqslant 0,

\]



сходяшийся во всей плоскости $\mathbb{C}, \lim _{k \rightarrow \infty} \sqrt[k]{\left|a_{k}\right|}=0$.



Простейшая Ц. Ф.- целая рациональная функ ция (или многочлен, полином) - функция вида



\[

w=P_{n}(z)=a_{0} z^{n}+a_{1} z^{n-1}+\ldots+a_{n}

\]



где $n$ - целое неотрицательное, коэффициенты $a_{0}, \ldots, a_{n}$ действительные или комплексные числа, $z$ - действительное или комплексное переменное.



ЧЕМПЛЕНА TEOPEМА - см. Ударные волны.



ЦЕНТР - один ИЗ вИдов особых точек дифференциального уравнения. В окрестности Этой точкИ все Интегральные кривые являются замкнутыми и содержат ее внутри себя (см. рис.). Ц. принадлежит к числу таких особых точек, характер к-рых, вообше говоря, не сохраняется при малых изменениях правой части уравнения. ЧЕНTP г р у п п - множество всех элементов $z$ груп-



\begin{center}

\includegraphics[max width=\textwidth]{2023_07_08_ae95fd759310c41356c9g-1}

\end{center}



Центр. пы $G$, удовлетворяюших условию $z g=g z$ для всех $g$ из $G$. Ц.- нормальный делитель в $G$. Ю. А. Неретин. ЦЕНТР ИНЕРЦИИ, ц е н т р м а с с,- геометрическая точка, положение которой характеризует распределение масс в теле или механической системе. Координаты Ц. и. определяются формулами:



\[

x_{c}=\sum m_{k} x_{k} / M, \quad y_{c}=\sum m_{k} y_{k} / M, \quad z_{c}=\sum m_{k} z_{k} / M

\]



или для тела при непрерывном распределении масс:



\[

x_{c}=\frac{1}{M} \int_{(V)} \rho x d V, \quad y_{c}=\frac{1}{M} \int_{(V)} \rho y d V, \quad z_{c}=\frac{1}{M} \int_{(V)} \rho z d V,

\]



где $m_{k}-$ массы материальных точек, образуюших систему, $x_{k}, y_{k}, z_{k}-$ координаты этих точек, $\boldsymbol{M}=\sum \boldsymbol{m}_{\boldsymbol{k}}-$ масса системы, $\rho(x, y, z)$ - плотность, $V$ - объем. Понятие о Ц. и. отличается от понятия о центре тяхести. тем, что последнее имеет смысл только для твердого тела, находяшегося в одно-





\end{document}